% TOP_PROBLEM: {"id": 1461, "title": "Vaught Conjecture", "category_id": 4, "category": {"id": 4, "name": "algebra", "display_name": "Algebra", "description": "Group theory, ring theory, field theory, and algebraic structures.", "slug": "algebra", "order_index": 4, "created_at": "2026-07-31T15:26:25.671Z"}, "status": "open"}
% FIRST_POSTED: null
% ENTRY_KIND: statement_only
% Imported from: batch07.tex; UnsolvedMath ID 1461
\documentclass[11pt]{article}
\usepackage[margin=25mm]{geometry}
\usepackage{fontspec}
\setmainfont{Latin Modern Roman}
\usepackage{amsmath,amssymb,mathtools}
\usepackage{unicode-math}
\setmathfont{Latin Modern Math}
\usepackage{xurl,hyperref}
\hypersetup{colorlinks=true,urlcolor=blue}
\setcounter{secnumdepth}{0}
\setlength{\parindent}{0pt}
\setlength{\parskip}{5pt}
\setlength{\emergencystretch}{3em}
\newcommand{\sref}[2]{\hyperlink{source:#1:#2}{[S#2]}}
\newcommand{\cataloguescope}[1]{\paragraph{Recorded target.}#1}
\begin{document}
% UnsolvedMath ID 1461; source catalogue code MOD-004
\setcounter{section}{1460}
\section{Vaught Conjecture}\label{1461}
{\small\sffamily UnsolvedMath ID 1461 · Algebra}
\subsection{Definitions and mathematical statement}

\paragraph{Shared notation.}$\mathbb N=\{1,2,\ldots\}$, and $\mathbb Z,\mathbb Q,\mathbb R,\mathbb C$ denote the integers, rationals, reals and complex numbers. Any other conventions are specified locally.


Let $L$ be a countable first-order language and $T$ a complete consistent theory in $L$, meaning that for each $L$-sentence exactly one of it and its negation is a consequence of $T$. Count models up to $L$-isomorphism, not as differently labeled structures. Let $I(T,\aleph_0)$ be the number of isomorphism classes of models of $T$ whose domains are countably infinite. Here $\aleph_0=|\mathbb N|$ and $2^{\aleph_0}=|\mathcal P(\mathbb N)|$ is the cardinality of the continuum. If $T$ has only finite models, then $I(T,\aleph_0)=0$.
\paragraph{Statement recorded in the supplied source.}
Does every complete consistent theory in a countable first-order language satisfy
\[
 I(T,\aleph_0)\le\aleph_0\qquad\text{or}\qquad I(T,\aleph_0)=2^{\aleph_0}?
\]
Thus the number of nonisomorphic countably infinite models must be finite, countably infinite, or continuum, without an intermediate uncountable value. The scope is finitary first-order theories, not all infinitary sentences or the stronger topological Vaught conjecture. \sref{1461}{2}
\subsection{Short English statement}
Must the number of nonisomorphic countable models of any complete countable-language theory be either at most countable or equal to the continuum?
\subsection{Sources}
\begin{description}
\item[\hypertarget{source:1461:1}{S1}] ULAM.AI and the UnsolvedMath Contributors, UnsolvedMath: A Curated Collection of Open Mathematics Problems, record 1461 (MOD-004). CC BY 4.0. \url{https://huggingface.co/datasets/ulamai/UnsolvedMath}.
\item[\hypertarget{source:1461:2}{S2}] Antonio Montalbán, A Computability Theoretic Equivalent to Vaught’s Conjecture, arXiv:1206.5682v2 (2013), Section 1, first paragraph and distinction between first-order and infinitary variants. \url{https://arxiv.org/pdf/1206.5682v2}.
\end{description}
\label{end:1461}
\end{document}
